\begin{abstract}\noindent
We revisit the tree-type neural network of the course note, in which a node's misclassified
samples are corrected by two child neurons and the children are grown recursively until the
training set is learned exactly, and replace every neuron by a small neural network. We give the
node problem as a neural version of the note's slack-minimisation LP, derive the closed-form solution
of the connection-weight LP, add a closed-form \emph{isolation node} that guarantees termination,
and generalise Table~II to $K$ classes in three ways. Nodes read features from frozen trunks
(pixels, a CNN trained from scratch, an ImageNet ResNet-18, or a per-node choice between trunks).
Every tree reaches \textbf{100\% training accuracy} on CIFAR-10 (50k images) and Cats vs Dogs (20k
images). The central finding is that \emph{how} a tree interpolates decides what interpolation costs:
learned neural correctors generalise their corrections and break two to three test predictions for
every one they fix (CIFAR-10: 84.8\%$\to$81.8\%), whereas closed-form isolation nodes interpolate
benignly (74.5\%$\to$74.6\%) but need one hidden unit per memorised image. With 20\% label noise the
root network shows a textbook double descent (test error 33\%$\to$54\%$\to$39\%, peak exactly where it
first interpolates); the tree interpolates at every width and lands on the over-parameterised branch
of the same curve. Receptive fields keep the trunk's spatial scale at every depth, but deep nodes memorise atypical,
cluttered or mislabelled photographs and increasingly recognise them by exclusion: their hidden units
veto the negatives and the positives are fired by the output bias alone. Decision-preserving pruning,
cost-complexity pruning and per-node trunk selection shrink the network substantially.
\end{abstract}
